\documentclass[11pt]{article}

\usepackage{amsmath,amssymb,amsthm}
\usepackage[margin=1.1in]{geometry}
\usepackage{bm}
\usepackage{booktabs}
\usepackage{array}
\usepackage[colorlinks=true,linkcolor=blue,citecolor=blue,urlcolor=blue]{hyperref}

\input{macros.inc}

\numberwithin{equation}{section}
\sloppy
\setlength{\emergencystretch}{3em}
\newtheorem{proposition}{Proposition}[section]
\newtheorem{lemma}[proposition]{Lemma}
\newtheorem{corollary}[proposition]{Corollary}
\theoremstyle{remark}
\newtheorem{remark}[proposition]{Remark}

% --- notation ---------------------------------------------------------------
\newcommand{\R}{\mathbb{R}}
\newcommand{\eps}{\varepsilon}
\newcommand{\dV}{\,\mathrm{d}V}
\newcommand{\dS}{\,\mathrm{d}S}
\newcommand{\tr}{\operatorname{tr}}
\newcommand{\sym}{\operatorname{sym}}
\newcommand{\dev}{\operatorname{dev}}
\newcommand{\vu}{\mathbf{u}}
\newcommand{\vv}{\mathbf{v}}
\newcommand{\vw}{\mathbf{w}}
\newcommand{\vn}{\hat{\mathbf{n}}}
\newcommand{\vt}{\mathbf{t}}
\newcommand{\vf}{\mathbf{f}}
\newcommand{\vx}{\mathbf{x}}
\newcommand{\vy}{\mathbf{y}}
\newcommand{\one}{\mathbf{1}}
\newcommand{\Ae}{A_{\eps}}
\newcommand{\Amu}{A_{\mu}}
\newcommand{\Mk}{M_{1/\kappa}}
\newcommand{\Mw}{M_{w}}
\newcommand{\Vh}{V_h}
\newcommand{\Qh}{Q_h}
\newcommand{\Pq}{\Pi_{Q}}
\newcommand{\code}[1]{\texttt{#1}}

\title{The pressure-split $(\mathbf{u}, p)$ solve for quasi-static
       (visco)elasticity:\\ formulation, discretisation and a
       parameter-robust block preconditioner}
\author{AdGIA method notes}
\date{}

\begin{document}
\maketitle

\noindent
This note documents an alternative linear solve for the displacement
problem of the quasi-static classes, selected with
\code{SetPressureSplitSolver()}. The displacement problem is rewritten
with the mean stress $p = \kappa\,\divg\vu$ as a second unknown,
discretised with a Taylor--Hood pair, and solved by MINRES with a
block-diagonal preconditioner whose quality does not depend on the ratio
of the bulk modulus to the (effective) shear modulus. In the language of
finite elements this is a \emph{mixed displacement--pressure
formulation}; in this codebase the word ``mixed'' is reserved for the
gravity-formulation axis (the \code{Mixed\ldots Problem} classes of
\code{doc/quasi\_static\_models.tex}), so the method is called \emph{the
pressure-split $(\vu, p)$ solve} throughout. The time layer whose
operators it solves is described in \code{doc/viscoelasticity.md}, the
problem classes in \code{doc/quasi\_static\_models.tex}, and the
measurements it rests on in \code{doc/benchmarks.tex} (``Box: viscosity
contrasts'') and \code{doc/planning/solvers.md} (``Near-incompressible
and high-contrast effective operators''). The sibling solver note is
\code{doc/gauge\_penalty\_iteration.tex}.

Notation follows \code{doc/quasi\_static\_models.tex}: $d \in \{2,3\}$
is the space dimension, $\eps(\vu) = \sym D\vu$ the linearised strain,
$\dev e = e - \frac1d (\tr e)\one$, and $A:B = A_{ij}B_{ij}$. The body
$\Omega$ has boundary $\partial\Omega = \Gamma_D \cup \Gamma_N$, with the
displacement prescribed on $\Gamma_D$ (of positive measure) and a traction
$\vt$ on $\Gamma_N$; $\vn$ is the outward unit normal. $V = \{\vv \in
H^1(\Omega)^d : \vv|_{\Gamma_D} = 0\}$, $(\cdot,\cdot)$ and
$\norm{\cdot}$ are the $L^2(\Omega)$ inner product and norm.

\tableofcontents

%==========================================================================
\section{The per-solve operator and its two ill-conditioned regimes}
\label{sec:setting}
%==========================================================================

\subsection{The effective operator of viscoelastic stepping}
\label{sec:beta}

The isotropic generalised Maxwell body of \code{doc/viscoelasticity.md}
\S1 has
\begin{equation}
  \sigma = \kappa\,(\tr\eps)\,\one + 2\mu_\infty\,\mathsf d
           + \sum_k 2\mu_k\,(\mathsf d - m_k),
  \qquad
  \dot m_k = \frac{\mathsf d - m_k}{\tau_k},
  \qquad \mathsf d := \dev\eps(\vu),
  \label{eq:maxwell}
\end{equation}
with trace-free internal variables $m_k$, branch shear moduli $\mu_k$ and
relaxation times $\tau_k(\vx)$; the bulk modulus does not relax. The
schemes that take one elastic solve per step (or stage) with an
\emph{effective} operator eliminate $m_k^{n+1}$ from the stress at the new
time. Write $h_k = \Delta t/\tau_k$ and $\mathsf d^n = \dev\eps(\vu^n)$.

\paragraph{Backward Euler.} The update $(m_k^{n+1} - m_k^n)/\Delta t =
(\mathsf d^{n+1} - m_k^{n+1})/\tau_k$ gives $m_k^{n+1} = (m_k^n + h_k
\mathsf d^{n+1})/(1+h_k)$, hence $\mathsf d^{n+1} - m_k^{n+1} =
(\mathsf d^{n+1} - m_k^n)/(1+h_k)$ and
\begin{equation}
  \sigma^{n+1} = \kappa\,(\tr\eps^{n+1})\one
   + 2\Bigl(\mu_\infty + \sum_k \beta_k \mu_k\Bigr)\mathsf d^{n+1}
   - \sum_k 2\mu_k \beta_k\, m_k^n,
  \qquad
  \beta_k^{\mathrm{BE}} = \frac{1}{1+h_k}.
  \label{eq:beta-be}
\end{equation}
An SDIRK stage is a backward-Euler step of length $\gamma\Delta t$, so its
weights are $1/(1+\gamma h_k)$.

\paragraph{Exponential trapezoid.} The internal-variable equation is
integrated exactly over the step with $\mathsf d$ linear in time,
$\mathsf d(t_n + s) = (1 - s/\Delta t)\mathsf d^n + (s/\Delta t)\mathsf
d^{n+1}$:
\[
  m_k^{n+1} = e^{-h_k} m_k^n
   + \frac{1}{\tau_k}\int_0^{\Delta t} e^{-(\Delta t - s)/\tau_k}\,
     \mathsf d(t_n + s)\,\mathrm{d}s .
\]
With $r = (\Delta t - s)/\tau_k$, so that $s/\Delta t = 1 - r/h_k$, and
$\int_0^h r e^{-r}\,\mathrm dr = 1 - e^{-h} - h e^{-h}$,
\[
  m_k^{n+1} = e^{-h_k} m_k^n + a(h_k)\,\mathsf d^n + b(h_k)\,\mathsf d^{n+1},
  \qquad
  a(h) = \frac{1 - e^{-h}}{h} - e^{-h},
  \quad
  b(h) = 1 - \frac{1 - e^{-h}}{h},
\]
the weights of \code{ExponentialTrapezoidWeights}. Substituting into
\eqref{eq:maxwell} at $t_{n+1}$,
\begin{equation}
\begin{gathered}
  \sigma^{n+1} = \kappa\,(\tr\eps^{n+1})\one
   + 2\Bigl(\mu_\infty + \sum_k \beta_k \mu_k\Bigr)\mathsf d^{n+1}
   - \sum_k 2\mu_k \bigl(e^{-h_k} m_k^n + a(h_k)\,\mathsf d^n\bigr),
  \\
  \beta_k^{\mathrm{ET}} = 1 - b(h_k) = \frac{1 - e^{-h_k}}{h_k}.
\end{gathered}
  \label{eq:beta-et}
\end{equation}
ETD1 and the explicit Runge--Kutta stages solve with the unrelaxed
modulus ($\beta_k = 1$).

\paragraph{The per-solve problem.} In every case the new displacement
solves: find $\vu \in \vu_D + V$ with
\begin{equation}
  a(\vu,\vv) := \int_\Omega \kappa\,(\divg\vu)(\divg\vv)
   + 2\mu\,\dev\eps(\vu):\dev\eps(\vv)\dV = \ell(\vv)
  \quad \forall \vv\in V,
  \qquad
  \mu := \mu_\infty + \sum_k \beta_k \mu_k ,
  \label{eq:primal}
\end{equation}
where $\ell(\vv) = \int_{\Gamma_N}\vt\cdot\vv\dS + \langle \vf,\vv\rangle$
collects the boundary traction and the memory forces (the last sums in
\eqref{eq:beta-be}--\eqref{eq:beta-et}, deviatoric dual vectors
assembled by the viscoelastic operator). From \S\ref{sec:split} on, $\mu$
always denotes this \emph{effective} shear modulus of the solve; it
coincides with the unrelaxed $\mu_U = \mu_\infty + \sum_k\mu_k$ when no
weights are set. The weights are nodal fields (the $\tau_k$ vary in space)
and $\mu$ is evaluated as the coefficient chain $\mu_\infty + \sum_k
\beta_k\mu_k$ at quadrature points (\code{doc/viscoelasticity.md} \S4).
The form \eqref{eq:primal} equals $\lambda(\divg\vu)(\divg\vv) +
2\mu\,\eps(\vu):\eps(\vv)$ with $\lambda = \kappa - 2\mu/d$, by the
orthogonal split $|e|^2 = |\dev e|^2 + \frac1d(\tr e)^2$.

\paragraph{Properties of the weights.} Both $\beta^{\mathrm{BE}}(h) =
1/(1+h)$ and $\beta^{\mathrm{ET}}(h) = (1-e^{-h})/h$ decrease
monotonically from $1$ at $h = 0$ to $0$ as $h\to\infty$, with
$\beta^{\mathrm{BE}} = 1 - h + O(h^2)$, $\beta^{\mathrm{ET}} = 1 - h/2 +
O(h^2)$ for small $h$, and
\begin{equation}
  h\,\beta(h) \to 1 \quad (h\to\infty),
  \qquad\text{so}\qquad
  \mu_k\beta_k \;\longrightarrow\; \frac{\mu_k\tau_k}{\Delta t}
  = \frac{\eta_k}{\Delta t},
  \label{eq:collapse}
\end{equation}
with $\eta_k = \mu_k\tau_k$ the branch viscosity. A branch that is relaxed
within the step contributes the shear stiffness of its \emph{viscosity
over the step length}, not of its elastic modulus.

\subsection{Two regimes with one root}
\label{sec:regimes}

The conditioning of \eqref{eq:primal} is governed pointwise by the ratio
\begin{equation}
  r(\vx) := \frac{\kappa(\vx)}{\mu(\vx)},
\end{equation}
and, secondarily, by the spatial contrast $\max\mu/\min\mu$.
\begin{itemize}
\item \emph{Near incompressibility.} $r$ is large because the material
  is: $\kappa/\mu_U \gg 1$. For the elastic Earth $\kappa/\mu_U$ is of
  order two, so this regime arises from the artificial and limiting
  constructions (the near-zero shear of fluid regions below) rather than
  from rock.
\item \emph{Effective-shear collapse.} With $\kappa$ unrelaxed and
  \eqref{eq:collapse}, a Maxwell region ($\mu_\infty = 0$) in which
  $\Delta t\gg\tau$ has
  \begin{equation}
    r \;\approx\; \frac{\kappa\,\Delta t}{\eta}
      = \frac{\kappa}{\mu_U}\, h,
  \end{equation}
  unbounded in $\Delta t$. The L-stable and exponential schemes are
  designed to stride over fast branches ($\Delta t\gg\tau$) with no loss of
  accuracy; the benchmarks confirm that their time accuracy does not
  depend on the viscosity contrast (\code{doc/benchmarks.tex}, ``Box:
  viscosity contrasts''). The price of striding is paid here, in the
  conditioning of the solve.
\item \emph{Multi-branch rheologies collapse in stages.} With several
  branches, $\mu(\Delta t;\vx) = \mu_\infty + \sum_k \mu_k\,
  \beta(\Delta t/\tau_k(\vx))$ is a staircase in $\Delta t$: branch $k$
  drops from $\mu_k$ to $\approx\eta_k/\Delta t$ around $\Delta t\approx
  \tau_k$, and the region where it has collapsed, $\{\vx :
  \tau_k(\vx)\ll\Delta t\}$, grows with $\Delta t$. With $\mu_\infty > 0$
  the collapse saturates at $r \le \kappa/\mu_\infty$; with $\mu_\infty =
  0$ it does not. For instance the \code{prony\_wide} body of the box
  benchmarks ($\mu_\infty = 0.1$, four branches $0.225$ at $\tau =
  10^{-3},10^{-1},10,10^{3}$) has, under the exponential trapezoid,
  $\mu = 1$ unrelaxed, $\mu \approx 0.56$ at $\Delta t = 1$ (two branches
  collapsed) and $\mu \approx 0.12$ at $\Delta t = 10^4$. Every change of
  step size is a different operator, so an adaptive stepper visits a
  continuum of them.
\end{itemize}
In the viscoelastic box benchmarks $\kappa/\mu_U = 5/3$; all their
ill-conditioning is collapse.

\subsection{Why algebraic multigrid on the primal operator degrades}
\label{sec:amg}

Split the energy as $a(\vv,\vv) = 2\mu\norm{\dev\eps(\vv)}^2 +
\kappa\norm{\divg\vv}^2$ (constant coefficients, for the argument). For
$r\gg1$ the operator has two very different scales: fields with an
$O(1)$ relative divergence cost $\sim\kappa$, fields with small divergence
cost $\sim\mu$. The \emph{near-null space} of $A$ relative to its scale
--- in algebraic multigrid (AMG) terms, the algebraically smooth vectors,
those with Rayleigh quotient small against the diagonal --- therefore
consists of the (nearly) divergence-free fields, $\norm{\divg\vv}^2\lesssim
r^{-1}\norm{\dev\eps(\vv)}^2$. In the continuum this is an
infinite-dimensional space with members at every length scale; discretely
it contains, for instance, the interpolants of smooth solenoidal fields at
every wavelength down to a few element sizes.

Two-grid theory requires (the weak approximation property) that each
algebraically smooth fine vector $\vv$ have a coarse representative
$P\vv_c$ with $\norm{\vv - P\vv_c}_{D}^2 \le K\,a(\vv,\vv)$, $D$ the
diagonal, and the two-grid contraction is then $1 - O(1/K)$. Nodal
coarsening with interpolation fitted to the translations (and, optionally,
rotations) builds a coarse space that represents smooth fields. It does
not represent nearly divergence-free fields oscillating below the coarse
scale, and, more fundamentally, interpolation does not commute with the
divergence: the coarse-grid interpolant of a nearly solenoidal field has
an $O(1)$ relative divergence, so its energy carries the factor $\kappa$
where the field's own carries $\mu$. The constant $K$ therefore grows like
$r$, and the coarse correction becomes ineffective on exactly the error
components the smoother leaves behind. Robust multigrid for this problem
in primal variables needs block smoothers whose patches span the local
divergence-free fields and prolongations that map coarse
divergence-free fields to fine ones \cite{Schoberl1999,LWXZ2007}; neither
is available in BoomerAMG \cite{HensonYang2002}. The measured cost is
consistent with this: a factor $2.5$ per solve from $r = 200$ to $2000$ on
a coarse ball at orders 2--3 (about $\sqrt{10}$, the growth of CG
iterations if the effective condition number grows like $r$), and on the
box ladder CG iterations rising from about 55 at contrast $C = 1$ to
150--340 on the layered cases at $C\ge10^4$ and 500--2500 on the graded
column at $C\ge10^6$, while the schemes that solve with the unrelaxed
operator stay at 50--65 (\code{doc/benchmarks.tex}). Supplying the rigid
rotations as interpolation vectors does not help, measured on the same
cases (\code{doc/planning/solvers.md}): they enrich the coarse space with
smooth fields, which it already represents, not with solenoidal ones.

%==========================================================================
\section{The pressure-split reformulation}
\label{sec:split}
%==========================================================================

\subsection{Stress split and weak saddle system}

Introduce the mean stress
\begin{equation}
  p := \kappa\,\divg\vu ,
  \qquad\text{so that}\qquad
  \sigma = p\,\one + 2\mu\,\dev\eps(\vu).
  \label{eq:pdef}
\end{equation}
($p$ is positive in extension: it is minus the thermodynamic pressure; the
code returns it under this sign from \code{Pressure()}.) Since $\one :
\eps(\vv) = \divg\vv$ and $\dev\eps(\vu):\eps(\vv) =
\dev\eps(\vu):\dev\eps(\vv)$, the virtual work is
$\int\sigma:\eps(\vv) = a_\mu(\vu,\vv) + b(\vv,p)$, and the constitutive
relation \eqref{eq:pdef} is imposed weakly with weight $1/\kappa$. The
problem becomes: find $(\vu,p)\in(\vu_D + V)\times L^2(\Omega)$ with
\begin{equation}
\begin{aligned}
  a_\mu(\vu,\vv) + b(\vv,p) &= \ell(\vv) && \forall\vv\in V,\\
  b(\vu,q) - c(p,q) &= 0 && \forall q\in L^2(\Omega),
\end{aligned}
\label{eq:saddle}
\end{equation}
\begin{equation}
  a_\mu(\vu,\vv) = \int_\Omega 2\mu\,\dev\eps(\vu):\dev\eps(\vv)\dV,
  \qquad
  b(\vv,q) = \int_\Omega q\,\divg\vv\dV,
  \qquad
  c(p,q) = \int_\Omega \frac{1}{\kappa}\,p\,q\dV .
  \label{eq:forms}
\end{equation}
In matrix form, with the discrete spaces of \S\ref{sec:disc},
\begin{equation}
  \mathcal A
  \begin{pmatrix} \vu \\ p \end{pmatrix}
  =
  \begin{pmatrix} \Amu & B^T \\ B & -\Mk \end{pmatrix}
  \begin{pmatrix} \vu \\ p \end{pmatrix}
  =
  \begin{pmatrix} \vf \\ 0 \end{pmatrix},
  \label{eq:block}
\end{equation}
symmetric and indefinite: $\Amu\succeq0$ (positive definite on the
discrete space with essential conditions, \S\ref{sec:norm}), $\Mk\succ0$.

\subsection{The natural boundary condition is unchanged}

Integrating the first equation of \eqref{eq:saddle} by parts with smooth
fields, $\int\sigma:\eps(\vv) = -\int(\divg\sigma)\cdot\vv +
\int_{\partial\Omega}(\sigma\vn)\cdot\vv$, so the strong form is
$-\divg\sigma = \vf$ in $\Omega$ and
\begin{equation}
  \bigl(p\,\one + 2\mu\,\dev\eps(\vu)\bigr)\vn = \vt
  \quad\text{on }\Gamma_N .
\end{equation}
By \eqref{eq:pdef} this is the physical traction condition of the primal
problem, $(\kappa\,\divg\vu\,\one + 2\mu\,\dev\eps(\vu))\vn = \vt$: the
same load functional $\ell$ serves both formulations, and the pressure
needs no boundary condition (it lives in $L^2$; in the fully clamped case
$\Gamma_N = \emptyset$ its mean is fixed by the $c$-term, \S\ref{sec:schur}).
This depends on using the symmetric-gradient form $a_\mu$; the Laplacian
form $\int\mu\nabla\vu:\nabla\vv$ common in Stokes codes would make the
natural condition a pseudo-traction.

\subsection{Exact elimination}
\label{sec:elim}

Assume $0<\kappa_0\le\kappa\le\kappa_1<\infty$. The second equation of
\eqref{eq:saddle} holds for every $q\in L^2$, and $c$ is an equivalent
inner product on $L^2$, so it is equivalent to $p = \kappa\,\divg\vu$
almost everywhere. Substituting into the first, $b(\vv,\kappa\,\divg\vu) =
\int\kappa(\divg\vu)(\divg\vv)$, and \eqref{eq:saddle} reduces to
\eqref{eq:primal}. Conversely, the primal solution with $p :=
\kappa\,\divg\vu$ solves \eqref{eq:saddle}. The two continuous problems are
equivalent for every finite $\kappa$; the saddle form remains meaningful
at $\kappa = \infty$ ($c = 0$, the incompressible problem), where the
primal form is not.

\subsection{The small-shear limit: penalised Stokes and the fluid region}
\label{sec:smallmu}

Scale the displacement by the shear modulus: in a region of constant
$(\kappa,\mu)$ put $\tilde\vu = 2\mu\,\vu$ and $t = 2\mu/\kappa$. Then
\eqref{eq:saddle} reads
\begin{equation}
  \int\dev\eps(\tilde\vu):\dev\eps(\vv) + \int p\,\divg\vv = \ell(\vv),
  \qquad
  \int q\,\divg\tilde\vu - t\int p\,q = 0,
\end{equation}
the \emph{penalised Stokes} system with ``velocity'' $\tilde\vu$ and
penalty parameter $t$. Near incompressibility ($\kappa\to\infty$) and
shear collapse ($\mu\to0$) are the same limit $t\to0$ of this one system:
the incompressible Stokes problem. In the physical variables the
pressure stays $O(\ell)$, the divergence $\divg\vu = p/\kappa$ stays
finite, and the solenoidal part of $\vu = \tilde\vu/(2\mu)$ is of size
$1/\mu$ --- driven only by the part of the load that is not a gradient.

\paragraph{Zero shear in a subregion.} Let $\mu = 0$ on an open region
$\Omega_f$ (a fluid) and $\mu>0$ on its complement $\Omega_s$, which is
connected and contains $\Gamma_D$ in its boundary. Then
\begin{equation}
  \ker\mathcal A = \bigl\{ (\vw,0) : \vw\in H^1_0(\Omega_f)^d,\
  \divg\vw = 0 \bigr\}
  \label{eq:kernel}
\end{equation}
($\vw$ extended by zero). Indeed, every such pair is in the kernel, since
$a_\mu(\vw,\cdot)$ only integrates over $\Omega_s$ where $\vw = 0$, and
$b(\vw,\cdot) = 0$. Conversely, testing $\mathcal A(\vu,p) = 0$ with
$(\vu,-p)$ gives $a_\mu(\vu,\vu) + c(p,p) = 0$, so $p = 0$ and
$\dev\eps(\vu) = 0$ on $\Omega_s$; a field with vanishing deviatoric
strain is a conformal Killing field (\S\ref{sec:norm}), and one vanishing on
$\Gamma_D$ vanishes on the connected $\Omega_s$, so $\vu\in
H^1_0(\Omega_f)^d$, and the second equation with $p = 0$ gives
$\divg\vu = 0$. The kernel is the welded fluid-relabelling space ---
the gauge subspace $\mathcal G$ of \code{doc/gauge\_penalty\_iteration.tex}.
A load is compatible with it iff $\ell(\vw) = 0$ for all such $\vw$, which
by de~Rham's theorem \cite[Ch.~I, Thm.~2.3 and Rem.~2.4]{GiraultRaviart}
means that $\ell$ restricted to $\Omega_f$ is a gradient, $\vf|_{\Omega_f}
= -\nabla\pi$: the hydrostatic compatibility condition.

\paragraph{The Maxwell fluid treatment.} The default fluid treatment of
the library (\code{SetFluid}, \code{doc/static\_fluid\_core.tex}) gives
the fluid an artificial Maxwell shear whose per-step effective modulus
$\eps\,\mu_c$ is driven towards zero by an escalating continuation. In
the pressure-split form that is exactly the limit just described: $p =
\kappa\,\divg\vu$ stays well defined in the fluid, the vanishing shear
appears in the scaled Schur weight $1/(2\eps\mu_c) + 1/\kappa$
(\S\ref{sec:prec}), and the residual indeterminacy of $\vu$ as $\eps\to0$ is
the relabelling kernel \eqref{eq:kernel}. That indeterminacy is physical
(and, with gravity and stratification, carries the band structure of
\code{doc/static\_fluid\_core.tex}); selecting the secular state within
it is the job of the Maxwell continuation, not of this solver. What the
pressure-split solve can remove is the $\kappa/(\eps\mu_c)$-dependence of
the cost of each continuation step.

%==========================================================================
\section{Discretisation}
\label{sec:disc}
%==========================================================================

\subsection{The Taylor--Hood pair}

The displacement space $\Vh\subset V$ is the library's continuous
Lagrange space of order $k\ge2$ (simplices or tensor-product elements);
the pressure space is the continuous Lagrange space one order below,
\begin{equation}
  \Qh = \{ q\in C^0(\bar\Omega) : q|_T\in\mathcal P_{k-1}(T)\
  (\text{resp. } \mathcal Q_{k-1}(T)) \},
\end{equation}
the Taylor--Hood pair \cite{TaylorHood}. For $k\ge2$ it satisfies the
uniform discrete inf-sup condition
\begin{equation}
  \inf_{0\ne q\in\Qh}\ \sup_{0\ne\vv\in\Vh}
  \frac{b(\vv,q)}{\norm{\nabla\vv}\,\norm{q}} \;\ge\; \beta > 0,
  \label{eq:infsup}
\end{equation}
with $\beta$ independent of $h$, on shape-regular meshes subject to mild
conditions near the boundary, in two and three dimensions and for both
element families \cite[\S8.8]{BBF}, \cite{Stenberg1984,Boffi1997}. When
$\Gamma_D = \partial\Omega$, the infimum runs over $\Qh\cap L^2_0(\Omega)$
(the constants are not seen by $b$); when $\Gamma_N\ne\emptyset$, over all
of $\Qh$. Since $|\sym G|\le|G|$, $\norm{\eps(\vv)}\le\norm{\nabla\vv}$
and \eqref{eq:infsup} holds a fortiori with $\norm{\eps(\vv)}$ in the
denominator, which is the form used below.

\subsection{Volumetric locking of the primal discretisation}
\label{sec:locking}

Discretise \eqref{eq:primal} directly in $\Vh$. With Korn's inequality
$\norm{\eps(\vv)}\ge c_K\norm{\nabla\vv}$ on $V$
\cite[\S11.2]{BrennerScott}, the energy norm $\norm{\vv}_a^2 = a(\vv,\vv)$
satisfies $2\mu c_K^2\norm{\nabla\vv}^2\le a(\vv,\vv)\le(2\mu +
d\kappa)\norm{\nabla\vv}^2$ (constant coefficients; the upper bound uses
$(\divg\vv)^2\le d\,|\eps(\vv)|^2$). Céa's lemma gives energy-optimality,
$\norm{\vu - \vu_h}_a = \min_{\vv\in\Vh}\norm{\vu - \vv}_a$, but the
$H^1$ quasi-optimality constant is the ratio of these bounds, $O(r)$.

The mechanism is the divergence constraint. The discrete energy is
bounded, $a(\vu_h,\vu_h) = \ell(\vu_h)\le C\norm{\ell}^2/\mu$, so
$\norm{\divg\vu_h}^2\le C/(\kappa\mu)\to0$ as $r\to\infty$, and $\vu_h$
tends to an element of the \emph{strongly} divergence-free subspace $Z_h =
\{\vv\in\Vh : \divg\vv = 0\}$. Whether that subspace is rich enough to
approximate the incompressible limit is a property of the element and the
mesh, and for low-order $C^0$ spaces it is not. The extreme case is linear
triangles on a fully clamped simply connected polygon: with $N_T$
triangles, $N_I$ interior and $N_B$ boundary vertices, Euler's formula and
the edge count $3N_T = 2N_E - N_B$ give $N_T = 2N_I + N_B - 2$. There are
$2N_I$ unknowns and $N_T$ elementwise constraints $\divg\vv|_T = 0$, of
which at most one is dependent ($\sum_T|T|\divg\vv|_T = 0$), so for
$N_B\ge3$ the constraint matrix has at least as many independent rows as
columns and generically $Z_h = \{0\}$: the discrete solution tends to zero
(total locking). For higher orders on general meshes $Z_h$ is nontrivial
but approximates poorly, and the error constant deteriorates with $r$
\cite{BabuskaSuri1992}, \cite[\S11.3]{BrennerScott}, \cite[Ch.~VI
\S3]{Braess}; only for sufficiently high order on non-degenerate meshes is
$Z_h$ rich enough \cite{ScottVogelius1985}.

\subsection{Why the split is locking-free}
\label{sec:reduced}

Eliminate the discrete pressure from \eqref{eq:block}: $p_h =
\Mk^{-1}B\vu_h$, and $\vu_h$ solves
\begin{equation}
  \bigl(\Amu + B^T\Mk^{-1}B\bigr)\vu_h = \vf ,
  \qquad
  \vv^TB^T\Mk^{-1}B\vv = \sup_{q\in\Qh}\frac{(\divg\vv,q)^2}{(q/\kappa,q)}
  = \kappa\,\norm{\Pq\divg\vv}^2
  \label{eq:reduced}
\end{equation}
(the last equality for constant $\kappa$; $\Pq$ is the $L^2$ projection
onto $\Qh$). The pressure-split discretisation is the primal one with the
volumetric energy $\kappa\norm{\divg\vv}^2$ replaced by
$\kappa\norm{\Pq\divg\vv}^2$: only the $\Qh$-component of the divergence
is penalised. This is the classical equivalence of mixed methods with
selective reduced integration \cite{MalkusHughes1978}. The constrained
subspace in the limit is now the \emph{weakly} divergence-free
$\{\vv\in\Vh : \Pq\divg\vv = 0\}$, which under \eqref{eq:infsup}
approximates the continuous divergence-free space as well as $\Vh$
approximates $H^1$ \cite[\S5.1]{BBF}. Uniform stability of
\eqref{eq:block} in norms that do not involve $\kappa$, and hence
$\kappa$-independent quasi-optimality of $(\vu_h,p_h)$, is
Corollary~\ref{cor:stability} below; the general theory of penalised
saddle-point problems is \cite[Ch.~III \S4]{Braess},
\cite[\S4.3]{BBF}.

Two consequences. First, the split and the primal discretisations are
different discretisations of the same problem, and differ at
discretisation-error level where the primal one is sound. Second,
\eqref{eq:reduced} is not a good way to \emph{solve}: the reduced operator
has the same $\kappa$-dependent near-null space as the primal one (and
$\Mk^{-1}$ is dense). The saddle form is kept, and the preconditioner
does the work.

\subsection{Pressure jumps at material interfaces}
\label{sec:jump}

Across an interface $\Sigma$ where $\mu$ jumps (on element faces, by the
house meshing rule), traction continuity $\jump{\sigma\vn} = 0$ gives
\begin{equation}
  \jump{p} = -\jump{2\mu\,\vn\cdot\dev\eps(\vu)\,\vn},
\end{equation}
generally nonzero: the mean stress is discontinuous at a shear-modulus
contrast. A continuous $\Qh$ cannot represent the jump: for a pressure with
jump $J$ along $\Sigma$, every $q\in\Qh$ differs from it by about $|J|/2$
on a strip of width $O(h)$ along $\Sigma$, so
$\inf_{q\in\Qh}\norm{p - q}\gtrsim|J|\,(h|\Sigma|)^{1/2}$. Through the
quasi-optimality of Corollary~\ref{cor:stability} the pressure, and with
it the displacement in the energy norm, then converge only at
$O(h^{1/2})$ near the interface. In the laterally uniform \emph{column}
of the box benchmarks this is visible in its simplest form: there $u =
u(z)$ is piecewise linear with kinks on element faces, the primal
discretisation is exact in space, and the split is not ($p = \kappa u'$
jumps at every layer interface). A pressure space continuous within each
material region and discontinuous across interfaces would represent the
jump; it is not what is implemented.

This is distinct from the $O(1)$ discrepancies seen at extreme contrast.
In the two-region test with a weak-branch weight $10^{-6}$ (an effective
$r\gtrsim10^6$ in one region) the primal and split solutions disagree at
$O(1)$ on coarse meshes. That is the primal discretisation locking
(\S\ref{sec:locking}), a correctness issue of the primal route, not an
error of the split; at moderate contrast the two agree to the
interface-smearing level (\S\ref{sec:verification}).

%==========================================================================
\section{The block-diagonal preconditioner}
\label{sec:prec}
%==========================================================================

\subsection{Operator preconditioning}

The framework is that of Mardal and Winther \cite{MardalWinther2011}: if
the block operator $\mathcal A:X\to X'$ of a symmetric problem is an
isomorphism with $\norm{\mathcal A}$ and $\norm{\mathcal A^{-1}}$ bounded
independently of the parameters, in some parameter-dependent norm on $X =
V\times Q$, then the Riesz map $\mathcal P$ of that norm (or anything
spectrally equivalent to it) is a parameter-robust preconditioner:
the eigenvalues of $\mathcal P^{-1}\mathcal A$ lie in a fixed set bounded
away from zero, and MINRES \cite{PaigeSaunders1975} converges at a
parameter-independent rate. For the present problem the norm is
\begin{equation}
  \norm{(\vv,q)}_X^2 = \norm{\sqrt{2\mu}\,\eps(\vv)}^2
   + \norm{\sqrt{w}\,q}^2,
  \qquad
  w := \frac{1}{2\mu} + \frac{1}{\kappa},
  \label{eq:Xnorm}
\end{equation}
with Riesz map $\mathcal P = \operatorname{diag}(\Ae, \Mw)$,
\begin{equation}
  \vv^T\Ae\vv = \int 2\mu\,|\eps(\vv)|^2,
  \qquad
  q^T\Mw q = \int\Bigl(\frac{1}{2\mu} + \frac{1}{\kappa}\Bigr)q^2 ,
\end{equation}
and the large modulus enters only through the bounded weight
$1/\kappa\le w$: the operator that the AMG sees never contains $\kappa$. The weight interpolates between the two limits of
\S\ref{sec:smallmu}: $w\approx1/(2\mu)$ for $r\gg1$, the Stokes scaling,
and $w\approx1/\kappa$ in the compressible limit. The analysis of this
section is for this exact $\mathcal P$; the implemented one replaces both
blocks by spectrally equivalent approximations (\S\ref{sec:inexact}). The
pattern --- block-diagonal preconditioning of the penalised saddle with
the pressure mass matrix weighted by the sum of the two compliances --- is
the one analysed for mixed elasticity by Klawonn \cite{Klawonn1998}, and
for Stokes by Silvester and Wathen \cite{SilvesterWathen1994} and Rusten
and Winther \cite{RustenWinther1992}.

\subsection{The displacement norm: full gradient, not deviatoric}
\label{sec:norm}

The $(1,1)$ block of $\mathcal A$ is the deviatoric $\Amu$, but the
displacement norm is the full symmetric gradient $\Ae$. Note
\begin{equation}
  \Ae = \Amu + \frac{2\mu}{d}\,D,
  \qquad \vv^TD\vv = \norm{\divg\vv}^2 ,
\end{equation}
so $\Ae$ is the primal operator \eqref{eq:primal} of a fictitious material
with bulk modulus $2\mu/d$, i.e.\ Lamé $\lambda = 0$, Poisson ratio $0$:
the best-conditioned isotropic elasticity operator for nodal AMG, with no
volumetric near-null space. $\Amu$, in contrast, is the elasticity
operator of zero bulk modulus ($\lambda = -2\mu/d$, Poisson ratio $-1$),
the edge of the stability range, and its kernel on $H^1$ is the space of
conformal Killing fields:

\paragraph{The deviatoric kernel.} $\dev\eps(\vv) = 0$ says that $\eps(\vv)$
is pointwise a multiple of $\one$. In 2-D the conditions are $\partial_1v_1
= \partial_2v_2$ and $\partial_2v_1 + \partial_1v_2 = 0$: the
Cauchy--Riemann equations for
\begin{equation}
  F(z) = v_1 + i\,v_2, \qquad z = x_1 + i\,x_2 ,
\end{equation}
so the kernel is the set of holomorphic $F$ with $F\in H^1$ ---
infinite-dimensional (translations $F = c$, rotations and dilations $F =
\alpha z$, and every $z^n$, $e^{\alpha z}$, \ldots). In 3-D the conformal
Killing fields form a 10-dimensional space of polynomials of degree $\le2$
(translations, rotations, dilation, the special conformal fields).

\paragraph{Consequences for coercivity.} Consider the hypothesis
\[
  a_\mu(\vv,\vv)\;\ge\;\alpha\,\vv^T\Ae\vv \qquad \forall\vv\in\Vh .
  \qquad\qquad\text{(H)}
\]
\begin{itemize}
\item \emph{Fully clamped, constant $\mu$: (H) holds with $\alpha =
  \frac12$.} For $\vv\in H^1_0$, two integrations by parts give
  $\int\partial_jv_i\,\partial_iv_j = \int(\divg\vv)^2$, hence
  $\norm{\eps(\vv)}^2 = \frac12\norm{\nabla\vv}^2 +
  \frac12\norm{\divg\vv}^2$ and $\norm{\nabla\vv}^2 =
  \norm{\nabla\times\vv}^2 + \norm{\divg\vv}^2\ge\norm{\divg\vv}^2$. So
  $\norm{\dev\eps(\vv)}^2 = \norm{\eps(\vv)}^2 - \frac1d\norm{\divg\vv}^2
  \ge\frac12\norm{\nabla\vv}^2\ge\frac12\norm{\eps(\vv)}^2$. For
  piecewise-constant $\mu$ this gives $\alpha\ge\frac12\mu_{\min}/\mu_{\max}$.
\item \emph{3-D, $\Gamma_D$ of positive area: (H) holds} with some
  $\alpha>0$ independent of $h$, by the generalised Korn inequality
  $\norm{\vv}_{H^1}\le C(\norm{\dev\eps(\vv)} + \norm{\vv})$ \cite{Dain2006}
  and a compactness argument, since no nonzero conformal Killing field
  vanishes on an open piece of a surface (the zero set of a nonzero
  infinitesimal Möbius motion has dimension at most one).
\item \emph{2-D with a traction boundary: (H) fails uniformly.} An
  inequality $\norm{\vv}_{H^1}\le C(\norm{\dev\eps(\vv)} + \norm{\vv})$
  would make the unit ball of the (infinite-dimensional) kernel compact
  in $H^1$ by Rellich's theorem, which is impossible. More concretely, on
  the box $(0,L)\times(0,H)$ clamped at $x_2 = 0$ (the box-benchmark
  shape) take $F_n(z) = e^{-nH}e^{-inz}$, of modulus $e^{n(x_2 - H)}$, and
  correct its boundary values with $W_n = e^{-nH}e^{-inx_1}(1 - x_2/H)$;
  then $\vu_n := (\operatorname{Re},\operatorname{Im})(F_n - W_n)\in V$,
  $\dev\eps(\vu_n) = -\dev\eps(W_n)$ and
  \[
    \norm{\dev\eps(\vu_n)}\le\norm{\nabla W_n}
      \le e^{-nH}(n + H^{-1})\sqrt{2LH},
    \qquad
    \norm{\nabla F_n}^2 = 2\!\int|F_n'|^2 = nL(1 - e^{-2nH}),
  \]
  so with Korn $\norm{\eps(\vu_n)}\ge c_K(\norm{\nabla F_n} -
  \norm{\nabla W_n})\sim c_K\sqrt{nL}$ the ratio
  $\norm{\dev\eps(\vu_n)}/\norm{\eps(\vu_n)}$ decays like $\sqrt n\,e^{-nH}$.
  Interpolating $\vu_n$ in $\Vh$ with $n\sim|\log h|$ shows that the
  discrete constant $\alpha_h$ of (H) tends to zero under refinement.
\end{itemize}
So an AMG for $\Amu$ would face, in 2-D, an infinite-dimensional
near-null space it cannot represent, and even an exact $\Amu^{-1}$ would
not define a uniform displacement norm. The saddle system as a whole is
nevertheless uniformly coercive in the $\Ae$-norm, because the
divergence that $\Amu$ does not see is controlled through $B$. The
precise statement needs one lemma.

\begin{lemma}[Korn-type inequality on the Taylor--Hood pair]
\label{lem:korn}
Let the mesh be affine simplicial and shape-regular and $k\ge2$. There is
$C_J$, depending only on $d$, $k$ and the shape regularity, such that
\begin{equation}
\begin{gathered}
  \norm{(I - \Pq)\divg\vv}\le C_J\norm{\dev\eps(\vv)}
  \qquad\text{and hence}\\
  \norm{\dev\eps(\vv)}^2 + \norm{\Pq\divg\vv}^2\ge\alpha_*\norm{\eps(\vv)}^2,
  \qquad \alpha_* = \frac{1}{1 + 2C_J^2/d},
\end{gathered}
  \label{eq:korn}
\end{equation}
for all $\vv\in\Vh$, with no condition on $\Gamma_D$.
\end{lemma}
\begin{proof}
On each element $\divg\vv|_T\in\mathcal P_{k-1}(T)$, so $\divg\vv$ is a
discontinuous piecewise polynomial of the degree of $\Qh$. Let $\mathcal
O$ be the averaging (Oswald) operator onto $\Qh$, which averages nodal
values at shared Lagrange nodes; for discontinuous piecewise polynomials
\cite[Thm.~2.2]{KarakashianPascal2003}
\[
  \norm{\phi - \mathcal O\phi}^2\le C_O\sum_{F\ \text{interior}}
  h_F\norm{\jump{\phi}}_{L^2(F)}^2 .
\]
Since $\Pq$ is the best approximation, $\norm{(I-\Pq)\divg\vv}\le
\norm{\divg\vv - \mathcal O\divg\vv}$. On an interior face $F$ with normal
$\vn$, $\vv$ is continuous, so its tangential derivatives are continuous
and $\jump{\nabla\vv} = \mathbf a\otimes\vn$ for some $\mathbf a$
(Hadamard). Write $\mathbf a = \alpha\vn + \beta\hat{\mathbf s}$ with
$\hat{\mathbf s}\perp\vn$ a unit vector. Then $\jump{\divg\vv} = \alpha$ and
$\jump{\dev\eps(\vv)} = \alpha(\vn\otimes\vn - \frac1d\one) +
\frac\beta2(\hat{\mathbf s}\otimes\vn + \vn\otimes\hat{\mathbf s})$, the two
terms being orthogonal, so
\[
  |\jump{\dev\eps(\vv)}|^2 = \Bigl(1 - \frac1d\Bigr)\alpha^2 +
  \frac{\beta^2}{2}\;\ge\;\frac{d-1}{d}\,|\jump{\divg\vv}|^2 .
\]
With the discrete trace inequality $h_F\norm{\phi}^2_{L^2(F)}\le
C_{\mathrm{tr}}\norm{\phi}^2_{L^2(T)}$ for polynomials on an element $T$
adjacent to $F$, and $d+1$ faces per element,
\[
  \sum_F h_F\norm{\jump{\divg\vv}}_F^2
  \le\frac{d}{d-1}\sum_F h_F\norm{\jump{\dev\eps(\vv)}}_F^2
  \le\frac{2d(d+1)C_{\mathrm{tr}}}{d-1}\norm{\dev\eps(\vv)}^2 ,
\]
which is the first inequality with $C_J^2 = 2d(d+1)C_OC_{\mathrm{tr}}/(d-1)$.
For the second, $\norm{\eps}^2 = \norm{\dev\eps}^2 + \frac1d\norm{\divg\vv}^2
\le\norm{\dev\eps}^2 + \frac2d(\norm{\Pq\divg\vv}^2 + C_J^2\norm{\dev\eps}^2)$,
and $2/d\le1$.
\end{proof}

The lemma says that a nearly conformal discrete field --- small
$\dev\eps$ --- has a divergence that is, up to $O(\norm{\dev\eps})$, in
$\Qh$ and so seen by the pressure test functions. (An exactly conformal
discrete field is a single holomorphic polynomial on all of $\Omega$, whose
divergence $2\operatorname{Re}F'$ is a continuous polynomial of degree
$k-1$, i.e.\ lies in $\Qh$.) On tensor-product or curved elements
$\divg\vv|_T$ is not in the local pressure space and the argument needs an
additional elementwise approximation term; it is not pursued here.

\subsection{The Schur complements}
\label{sec:schur}

There are two Schur complements: that of the operator,
$S_\mu = B\Amu^{-1}B^T + \Mk$, and that of the preconditioner's metric,
$S_\eps = B\Ae^{-1}B^T + \Mk$. For variable coefficients define the
weighted inf-sup constant
\begin{equation}
  \beta_\mu := \inf_{0\ne q\in\Qh}\ \sup_{0\ne\vv\in\Vh}
  \frac{b(\vv,q)}{\norm{\sqrt{2\mu}\,\eps(\vv)}\,
  \norm{q/\sqrt{2\mu}}} \;\ge\;\beta\,\sqrt{\mu_{\min}/\mu_{\max}},
\end{equation}
the bound following from \eqref{eq:infsup} with
$\norm{\sqrt{2\mu}\,\eps(\vv)}\le\sqrt{2\mu_{\max}}\norm{\nabla\vv}$ and
$\norm{q/\sqrt{2\mu}}\le\norm{q}/\sqrt{2\mu_{\min}}$. For constant $\mu$,
$\beta_\mu\ge\beta$; for two subdomains with piecewise-constant viscosity
the corresponding constant of the Stokes problem is bounded independently
of the viscosity ratio \cite{OlshanskiiReusken2006}, and for smoothly
varying viscosity the inverse-viscosity-weighted mass matrix is a good
Schur approximation in practice, with constants depending on the
variation \cite{GrinevichOlshanskii2009}.

\begin{proposition}[Schur complement equivalence]
\label{prop:schur}
Let $\kappa,\mu$ be bounded positive measurable functions (in particular,
piecewise constant), and $\Gamma_N\ne\emptyset$. Then
\begin{equation}
  \min(\beta_\mu^2,1)\,\Mw\;\preceq\;S_\eps\;\preceq\;d\,\Mw ,
  \label{eq:schur-eps}
\end{equation}
and for the operator's Schur complement
\begin{equation}
  \min(\beta_\mu^2,1)\,\Mw\;\preceq\;S_\mu\;\preceq\;\frac{d}{\alpha}\,\Mw
  \quad\text{under (H)}.
  \label{eq:schur-mu}
\end{equation}
\end{proposition}
\begin{proof}
\emph{Upper bound for $S_\eps$.} Pointwise $(\tr e)^2\le d\,|e|^2$, so
$|\divg\vv|\le\sqrt d\,|\eps(\vv)|$ and, by a weighted Cauchy--Schwarz
inequality,
\[
  b(\vv,q) = \int\frac{q}{\sqrt{2\mu}}\,\sqrt{2\mu}\,\divg\vv
  \le\sqrt d\,\norm{q/\sqrt{2\mu}}\,\norm{\sqrt{2\mu}\,\eps(\vv)} .
\]
Hence $q^TB\Ae^{-1}B^Tq = \sup_{\vv}b(\vv,q)^2/(\vv^T\Ae\vv)\le d\,q^TM_{1/(2\mu)}q$,
and $S_\eps\preceq dM_{1/(2\mu)} + \Mk\preceq d\,\Mw$ since $d\ge1$.
\emph{Lower bound.} By definition of $\beta_\mu$, $q^TB\Ae^{-1}B^Tq\ge
\beta_\mu^2q^TM_{1/(2\mu)}q$; adding $\Mk$, $S_\eps\succeq\min(\beta_\mu^2,1)
(M_{1/(2\mu)} + \Mk) = \min(\beta_\mu^2,1)\Mw$. The weights add pointwise,
so no condition on the coefficients is needed.
\emph{$S_\mu$.} Pointwise $|\dev e|\le|e|$, so $\Amu\preceq\Ae$, hence
$\Amu^{-1}\succeq\Ae^{-1}$ and $S_\mu\succeq S_\eps$: the lower bound
transfers. Under (H), $\Amu^{-1}\preceq\alpha^{-1}\Ae^{-1}$ and
$S_\mu\preceq\alpha^{-1}B\Ae^{-1}B^T + \Mk\preceq(d/\alpha)\Mw$.
\end{proof}

\begin{remark}[Fully clamped bodies]
\label{rem:clamped}
If $\Gamma_D = \partial\Omega$, then $b(\vv,1) = \int_{\partial\Omega}
\vv\cdot\vn = 0$ for all $\vv\in\Vh$: $B^T\one = 0$. For constant
coefficients $\one$ is then an exact eigenvector of the pencil
$(S_\eps,\Mw)$ (and of $(S_\mu,\Mw)$) with eigenvalue
$(1/\kappa)/w = 2\mu/(2\mu + \kappa)$, and $(0,\one)$ an exact eigenvector
of $\mathcal P^{-1}\mathcal A$ with eigenvalue $-2\mu/(2\mu+\kappa)$. Its
$\mathcal P$-orthogonal complement is invariant, consists of mean-free
pressures, and there the bounds hold with $\beta$ the inf-sup constant on
$\Qh\cap L^2_0$. So a fully clamped body has a single outlying eigenvalue,
tending to zero like $1/r$; it costs MINRES $O(\log r)$ additional
iterations in the worst-case bound (a factor $1 + |\lambda|_{\max}/|\lambda_o|$
for the polynomial that vanishes at the outlier) and is deflated in
practice. The constant pressure mode is fixed by the $c$-term alone,
$(p/\kappa,1) = \int_{\partial\Omega}\vu_D\cdot\vn$: it is the whole
pressure in the constant-strain verification gate.
\end{remark}

Proposition~\ref{prop:schur} shows that the true Schur complement $S_\mu$ is
equivalent to $\Mw$ when (H) holds, and that its upper bound is
lost exactly where (H) is (2-D with traction boundaries: $q =
\Pq\divg\vv$ for a nearly conformal $\vv$ has Rayleigh quotient
$\gtrsim1/\alpha_h$). The block preconditioner never needs that upper
bound. Its theory uses $S_\eps$, which is equivalent to $\Mw$
unconditionally, together with Lemma~\ref{lem:korn}.

\subsection{Spectrum of the preconditioned saddle operator}
\label{sec:spectrum}

To cover the implemented blocks, state the bounds for general symmetric
positive definite $P_u\approx\Ae$ and $P_p\approx\Mw$, with
\begin{equation}
  c_A P_u\preceq\Ae\preceq C_AP_u,
  \qquad
  c_S P_p\preceq\Mw\preceq C_SP_p
  \label{eq:equiv}
\end{equation}
($c_A = C_A = c_S = C_S = 1$ for the exact Riesz map).

\begin{proposition}[Eigenvalue inclusion]
\label{prop:spectrum}
Let $\kappa,\mu$ be constant with $\kappa\ge2\mu/d$ (non-negative Lamé
$\lambda$; the effective operators inherit this from the unrelaxed one,
since $\mu\le\mu_U$), let the mesh satisfy the hypotheses of
Lemma~\ref{lem:korn}, $\Gamma_N\ne\emptyset$, and $\gamma :=
\min(\beta^2,1)$. Every eigenvalue of $\mathcal P^{-1}\mathcal A$, $\mathcal
P = \operatorname{diag}(P_u,P_p)$, lies in $[-\nu_{\max},-\nu_{\min}]\cup
[\lambda_{\min},\lambda_{\max}]$ with
\begin{align}
  \lambda_{\max} &\le \tfrac12\Bigl(C_A + \sqrt{C_A^2 + 4dC_AC_S}\Bigr),
  &
  \nu_{\max} &\le \tfrac12\Bigl(C_S + \sqrt{C_S^2 + 4dC_AC_S}\Bigr),
  \label{eq:upper}\\
  \lambda_{\min} &\ge \min\bigl(1,\ \theta\,\alpha_*\,c_A\bigr),
  &
  \nu_{\min} &\ge \min\Bigl(\gamma',\ \tfrac12\bigl(-C_A +
  \sqrt{C_A^2 + 4\gamma'}\bigr)\Bigr),
  \label{eq:lower}
\end{align}
where $\theta = \min\bigl(1, c_S/((d+1)(1 + C_S))\bigr)$ and $\gamma' =
\min(c_A,1)\,c_S\,\gamma$. All bounds are independent of $\kappa$ (including
$\kappa = \infty$), of the scale of $\mu$, and of $h$.
\end{proposition}
\begin{proof}
Write $\tilde A = P_u^{-1/2}\Amu P_u^{-1/2}$, $\tilde B =
P_p^{-1/2}BP_u^{-1/2}$, $\tilde C = P_p^{-1/2}\Mk P_p^{-1/2}$. The eigenproblem
is $\tilde A\vx + \tilde B^T\vy = \lambda\vx$, $\tilde B\vx - \tilde C\vy =
\lambda\vy$. From \eqref{eq:equiv} and the pointwise inequalities used in
Proposition~\ref{prop:schur}: $0\preceq\tilde A\preceq C_A$ (as
$\Amu\preceq\Ae$); $\tilde B\tilde B^T\preceq dC_AC_S$ (as
$P_u^{-1}\preceq C_A\Ae^{-1}$ and $B\Ae^{-1}B^T\preceq dM_{1/(2\mu)}\preceq
d\Mw\preceq dC_SP_p$); $0\preceq\tilde C\preceq C_S$ (as $\Mk\preceq\Mw$);
and $\tilde B\tilde B^T + \tilde C\succeq\gamma'$ (as $P_u^{-1}\succeq
c_A\Ae^{-1}$, $P_p^{-1}\succeq c_S\Mw^{-1}$ and $S_\eps\succeq\gamma\Mw$).

\emph{Positive eigenvalues.} For $\lambda>0$, $\vx\ne0$ (if $\vx = 0$ then
$-\tilde C\vy = \lambda\vy$ forces $\lambda\le0$). Eliminating $\vy =
(\lambda + \tilde C)^{-1}\tilde B\vx$ and taking $|\vx| = 1$,
\begin{equation}
  \lambda = \vx^T\tilde A\vx + \vx^T\tilde B^T(\lambda + \tilde C)^{-1}
  \tilde B\vx .
  \label{eq:pos}
\end{equation}
Upper: $\lambda\le C_A + |\tilde B\vx|^2/\lambda\le C_A + dC_AC_S/\lambda$,
a quadratic inequality giving \eqref{eq:upper}. Lower: suppose $\lambda<1$.
Since $\tilde C\preceq C_S$, \eqref{eq:pos} gives $\lambda\ge\vx^T\tilde
A\vx + |\tilde B\vx|^2/(1 + C_S)$. In the original variables, with $\vv =
P_u^{-1/2}\vx$,
\[
  |\tilde B\vx|^2 = \vv^TB^TP_p^{-1}B\vv\ge c_S\,\vv^TB^T\Mw^{-1}B\vv
  = \frac{c_S}{w}\norm{\Pq\divg\vv}^2
  \ge\frac{2\mu\,c_S}{d+1}\norm{\Pq\divg\vv}^2 ,
\]
where $1/w = 2\mu\kappa/(2\mu + \kappa)\ge2\mu/(d+1)$ for
$\kappa\ge2\mu/d$. Hence, by Lemma~\ref{lem:korn},
\begin{align*}
  \lambda&\ge2\mu\norm{\dev\eps(\vv)}^2 + \frac{2\mu c_S}{(d+1)(1+C_S)}
  \norm{\Pq\divg\vv}^2\\
  &\ge\theta\,2\mu\bigl(\norm{\dev\eps(\vv)}^2 +
  \norm{\Pq\divg\vv}^2\bigr)\ge\theta\alpha_*\,\vv^T\Ae\vv
  \ge\theta\alpha_*c_A ,
\end{align*}
using $\vv^T\Ae\vv\ge c_A\vv^TP_u\vv = c_A$.

\emph{Negative eigenvalues} $\lambda = -\nu$, $\nu>0$. Now $\vy\ne0$, and
eliminating $\vx = -(\tilde A + \nu)^{-1}\tilde B^T\vy$ with $|\vy| = 1$,
\begin{equation}
  \nu = \vy^T\tilde B(\tilde A + \nu)^{-1}\tilde B^T\vy + \vy^T\tilde C\vy .
  \label{eq:neg}
\end{equation}
Upper: $(\tilde A + \nu)^{-1}\preceq\nu^{-1}$, so $\nu\le dC_AC_S/\nu + C_S$,
giving \eqref{eq:upper}. Lower: $(\tilde A + \nu)^{-1}\succeq(C_A +
\nu)^{-1}$, so with $s = |\tilde B^T\vy|^2$ and $c = \vy^T\tilde C\vy$,
$\nu\ge s/(C_A + \nu) + c\ge\min\bigl(1,(C_A + \nu)^{-1}\bigr)(s + c)$ and
$s + c\ge\gamma'$. Either $C_A + \nu\le1$ and $\nu\ge\gamma'$, or
$\nu(C_A + \nu)\ge\gamma'$; this is \eqref{eq:lower}.
\end{proof}

For the exact Riesz map the bounds are $\lambda_{\max},\nu_{\max}\le
\frac12(1 + \sqrt{1 + 4d})$ ($\approx2.30$ in 3-D), $\nu_{\min}\ge
\frac12(-1 + \sqrt{1 + 4\gamma})$ and $\lambda_{\min}\ge\min(1,
\alpha_*/(2(d+1)))$. The hypothesis (H) does not appear: the
positive lower bound comes from the lemma, in which the $B$-coupling
supplies the divergence that $\Amu$ does not see. This is the precise sense
in which the full-gradient operator, not the deviatoric one, is the right
displacement block of the preconditioner.

\begin{remark}
With an exact Schur complement in the preconditioner ($\mathcal P =
\operatorname{diag}(\Amu,S_\mu)$, $\Mk = 0$), $\mathcal P^{-1}\mathcal A$
would have only the three eigenvalues $1$, $\frac12(1\pm\sqrt5)$
\cite{MurphyGolubWathen2000}. That choice is unavailable here twice over:
$S_\mu$ is dense, and $\Amu$ is not a usable norm (\S\ref{sec:norm}).
\end{remark}

\subsection{Consequences: stability, locking-freeness, MINRES}

\begin{corollary}[Uniform discrete stability]
\label{cor:stability}
Under the hypotheses of Proposition~\ref{prop:spectrum} with the exact
Riesz map, the bilinear form $\mathfrak B((\vu,p),(\vv,q)) = a_\mu(\vu,\vv) +
b(\vv,p) + b(\vu,q) - c(p,q)$ satisfies, on $\Vh\times\Qh$ with the norm
\eqref{eq:Xnorm}, the inf-sup condition with constant $\gamma_h =
\min(\lambda_{\min},\nu_{\min})$ and is bounded with constant
$\norm{\mathfrak B}\le\max(\lambda_{\max},\nu_{\max})$ (the latter on all
of $V\times L^2$). Hence \cite{Babuska1971}
\begin{equation}
  \norm{(\vu - \vu_h, p - p_h)}_X\le
  \Bigl(1 + \frac{\norm{\mathfrak B}}{\gamma_h}\Bigr)
  \inf_{(\vv,q)\in\Vh\times\Qh}\norm{(\vu - \vv, p - q)}_X ,
\end{equation}
with a constant independent of $\kappa$ and $h$: the pressure-split
discretisation is free of volumetric locking.
\end{corollary}
\begin{proof}
$\mathcal P^{-1}\mathcal A$ is self-adjoint in the $\mathcal P$-inner
product, which is the $X$ inner product; for such an operator the inf-sup
constant and the norm in $X$ are the smallest and largest moduli of its
eigenvalues. Boundedness on the continuous space uses only the pointwise
inequalities $|\dev e|\le|e|$, $|\tr e|\le\sqrt d|e|$ and $1/\kappa\le w$.
\end{proof}

\paragraph{MINRES.} Preconditioned MINRES minimises the residual in the
$\mathcal P^{-1}$-norm over the Krylov space. If the spectrum of $\mathcal
P^{-1}\mathcal A$ lies in $[-a,-b]\cup[c,e]$ with $a-b = e-c$ (enlarge the
shorter interval otherwise), then \cite[Ch.~4]{ESW}, \cite[\S3.1]{Greenbaum1997}
\begin{equation}
  \frac{\norm{r_k}_{\mathcal P^{-1}}}{\norm{r_0}_{\mathcal P^{-1}}}\le
  2\left(\frac{\sqrt{ae} - \sqrt{bc}}{\sqrt{ae} + \sqrt{bc}}\right)^{\lfloor k/2\rfloor}.
  \label{eq:minres}
\end{equation}
For symmetric intervals ($a = e$, $b = c$) with ratio $\varkappa = e/c$ the
factor per two iterations is $(\varkappa - 1)/(\varkappa + 1)$, so a
relative tolerance $\tau$ needs about $\varkappa\ln(2/\tau)$ iterations.
By Proposition~\ref{prop:spectrum} $\varkappa$ is bounded independently of
$\kappa$, $\mu$ and $h$: the iteration count is \emph{parameter-robust}.
Note what the bound does not promise: equality with CG on a
well-conditioned primal problem. CG on a spectrum of condition number $K$
needs about $\frac12\sqrt K\ln(2/\tau)$ iterations; an indefinite spectrum
of interval ratio $\varkappa$ behaves like $K\approx4\varkappa^2$. Together
with the extra work per iteration (the $B$, $B^T$ and $\Mk$ products and a
second preconditioner block) this is the saddle overhead, measured at
about $2\times$ when nothing is ill-conditioned (\S\ref{sec:results}).

\subsection{What the theory does not cover}
\label{sec:inexact}

The bounds of Proposition~\ref{prop:spectrum} are uniform for constant
coefficients and exact Riesz maps; the implemented solve departs from that
in three ways, and they account for the residual dependences that are
measured.
\begin{itemize}
\item \emph{The displacement block.} $P_u^{-1}$ is one BoomerAMG V-cycle
  on $\Ae$. Because $\Ae$ has Poisson ratio $0$ and no $\kappa$, the AMG
  faces no volumetric near-null space; but with a variable $\mu$ it is a
  heterogeneous-coefficient elasticity operator, and $c_A$, $C_A$ depend
  on the coefficient's variation and, through it, on $h$. On the graded
  column the effective modulus spans five to six decades across the
  layer, smoothly.
\item \emph{The pressure block.} $P_p^{-1}$ is one Jacobi sweep on $\Mw$.
  The Jacobi-preconditioned mass matrix has eigenvalues within elementwise
  bounds that depend only on the element type for affine elements and
  constant weight \cite{Wathen1987} (e.g.\ $[\frac12,2]$ for linear
  triangles, $[\frac12,\frac52]$ for linear tetrahedra, $[\frac14,\frac94]$
  for bilinear rectangles); with a weight varying inside elements and with
  distorted or curved elements the constants $c_S$, $C_S$ also depend on
  the in-element variation of $w$ and of the Jacobian.
\item \emph{Variable coefficients in the saddle itself.} For variable
  $\mu$, $\beta$ is replaced by $\beta_\mu$ (contrast-uniform only in the
  settings cited in \S\ref{sec:schur}) and the lemma needs a weighted
  analogue whose contrast-uniformity is not established here. In 2-D with
  traction boundaries the constant $\alpha_*$ of the lemma is still
  uniform, but the true Schur complement is not equivalent to $\Mw$
  (\S\ref{sec:schur}); this does not enter the bounds.
\end{itemize}
None of these brings $\kappa$ back: they are dependences on the shear
coefficient's variation, on element shape and on $h$, which is what the
graded measurements show (iteration counts flat in the contrast ladder,
growing under refinement of the graded column).

%==========================================================================
\section{Boundary conditions and solve mechanics}
\label{sec:mechanics}
%==========================================================================

\paragraph{Essential conditions.} Let the displacement true dofs split
into free (0) and essential (e) ones, with prescribed values
$\vu_e = \vu_{\mathrm{bdr}}$. The full system \eqref{eq:block} reads
\[
  \begin{pmatrix}
    A_{00} & A_{0e} & B_0^T\\
    A_{e0} & A_{ee} & B_e^T\\
    B_0 & B_e & -\Mk
  \end{pmatrix}
  \begin{pmatrix}\vu_0\\ \vu_e\\ p\end{pmatrix}
  =
  \begin{pmatrix}\vf_0\\ \ast\\ 0\end{pmatrix} .
\]
The eliminated system replaces the essential rows by the identity
(\code{DIAG\_ONE}, as on the primal path) and moves the essential columns
to the right-hand side:
\begin{equation}
  \begin{pmatrix}
    A_{00} & 0 & B_0^T\\
    0 & I & 0\\
    B_0 & 0 & -\Mk
  \end{pmatrix}
  \begin{pmatrix}\vu_0\\ \vu_e\\ p\end{pmatrix}
  =
  \begin{pmatrix}\vf_0 - A_{0e}\vu_{\mathrm{bdr}}\\
  \vu_{\mathrm{bdr}}\\ -B_e\vu_{\mathrm{bdr}}\end{pmatrix} .
  \label{eq:eliminated}
\end{equation}
It is symmetric, so MINRES applies. On the $\vu$-row this is
\code{FormSystemMatrix}/\code{FormLinearSystem} on the deviatoric form.
For $B$ it is a \emph{rectangular} elimination of the essential trial
columns only (the pressure has no essential dofs): the assembled matrix is
$[B_0\ 0]$, and $B^T$ in the block operator is its transpose, whose
essential rows are therefore zero. The pressure row's right-hand side is
the boundary lift $-B_e\vu_{\mathrm{bdr}}$, produced by
\code{FormRectangularLinearSystem} applied to the grid function carrying
the boundary values with a zero right-hand side. The preconditioner's
$\Ae$ is eliminated on the same essential dofs, but its essential
diagonal is \emph{kept} at the assembled value rather than set to one:
with a collapsed effective shear the clamped boundary can sit in a
region where the $\kappa$-free block's entries are many orders below
one, and unit essential rows put an artificial jump there that the
AMG's coarsening mishandles. A preconditioner only needs to stay
symmetric positive definite, and the kept (assembled, positive)
diagonal is at the scale of its neighbours. The operator blocks keep
the unit diagonal, which the reduced system's consistency requires.

\paragraph{Warm starts.} The solve is warm-started in both unknowns: the
displacement from the problem's grid function (whose interior is copied
into the initial vector by \code{FormLinearSystem} with
\code{copy\_interior}) and the pressure from the previous solve's pressure
field. A relative tolerance in MINRES is relative to the initial
(preconditioned) residual, which for a good warm start is already tiny;
so, exactly as on the primal CG path, the absolute tolerance is set per
solve to the target a cold start would have had,
\begin{equation}
  \tau_{\mathrm{abs}} = \tau_{\mathrm{rel}}\,
  \bigl(\mathbf b^T\mathcal P^{-1}\mathbf b\bigr)^{1/2},
\end{equation}
with $\mathbf b$ the eliminated right-hand side and $\mathcal P^{-1}$ the
applied block preconditioner (\code{SetWarmStartTolerance}). This is the
norm in which MINRES measures its residual, so the semantics coincide
with the primal path's: stop when the preconditioned residual falls below
$\tau_{\mathrm{rel}}$ times the preconditioned norm of the data. When
$\mathbf b^T\mathcal P^{-1}\mathbf b\le0$ (zero data) the solve is
skipped and the solution set to zero.

\paragraph{Recovery.} After the solve the displacement block is copied
back into the vector that \code{FormLinearSystem} returned \emph{before}
\code{RecoverFEMSolution} is called (in serial that vector aliases the
grid function's storage and the recovery assumes the solver wrote through
it; \code{doc/mfem\_notes.md}), and the pressure grid function is set from
its true dofs; it is the warm start of the next solve and is exposed by
\code{Pressure()}.

%==========================================================================
\section{Implementation map}
\label{sec:impl}
%==========================================================================

\paragraph{Entry points} (\code{include/AdGIA/}\allowbreak
\code{quasi\_static\_problem.hpp}, \code{src/}\allowbreak
\code{quasi\_static\_problem.cpp}).
\code{SetPressureSplitSolver()} switches the route on (idempotent),
\code{ClearPressureSplitSolver()} returns to the primal solve and releases
every block, \code{UsingPressureSplit()} reports, and \code{Pressure()}
returns the mean-stress field on the internal pressure space.
\code{Solve()} dispatches to \code{PressureSplitSolve()} when the route is
on; that routine reassembles when the operator is dirty or its block
version lags the problem's operator version (a primal assembly made for
diagnostics through \code{SystemMatrix()} clears the dirty flag without
touching the blocks, so the split keys on the version), builds the
eliminated block system \eqref{eq:eliminated}, and runs MINRES.

\paragraph{Assembly} (\code{AssemblePressureSplitOperator()}). Five forms:
\begin{center}
\small
\begin{tabular}{@{}lll@{}}
\toprule
block & form & integrator \\
\midrule
$\Amu$ & $2\mu\,\dev\eps:\dev\eps$ & \code{ElasticityIntegrator(mu, -2/d, 1)} \\
$B$ & $(q,\divg\vu)$ & \code{VectorDivergenceIntegrator} (mixed form) \\
$\Mk$ & $(p/\kappa,q)$ & \code{MassIntegrator(1/kappa)}, weight $-1$ in the block \\
$\Ae$ (prec.) & $2\mu\,\eps:\eps$ & \code{ElasticityIntegrator(mu, 0, 1)} \\
$\Mw$ (prec.) & $((1/(2\mu) + 1/\kappa)p,q)$ & \code{MassIntegrator}, weight a sum of ratio coefficients \\
\bottomrule
\end{tabular}
\end{center}
(MFEM's \code{ElasticityIntegrator(c, q\_l, q\_m)} has Lamé parameters
$q_l c$, $q_m c$, so $\lambda = -2\mu/d$ gives the deviatoric form and
$\lambda = 0$ the full-gradient one.) In parallel the displacement block
of the preconditioner is \code{HypreBoomerAMG} on $\Ae$ with the house
nodal-systems settings of \code{SetupDefaultPreconditioner} (systems
options for the \code{byNODES} ordering, nodal coarsening, symmetric
$\ell_1$-Gauss--Seidel relaxation, so that the V-cycle is symmetric
positive definite as MINRES requires), and the pressure block is
\code{HypreDiagScale} on $\Mw$. In serial both blocks are Jacobi
(\code{DSmoother}): MINRES needs a symmetric positive definite
preconditioner, which rules out the serial Gauss--Seidel of the primal
path; the serial route is kept correct, not tuned. MINRES runs at the
problem's relative tolerance with \code{iterative\_mode} on.

\paragraph{The rheology side.} The route needs the pair $(\kappa,\mu)$ as
coefficients, supplied by \code{ElasticStiffness::IsotropicSplit(kappa,
mu)} (\code{include/AdGIA/rheology.hpp}, \code{src/rheology.cpp}). It
returns $\mu$ as the \emph{live} effective shear modulus: for an
isotropic elastic rheology the fixed pair; for an isotropic Maxwell
rheology the redirectable coefficient that the stiffness integrators
themselves read, which points at $\mu_U$ or, while relaxation weights are
set, at the chain $\mu_\infty + \sum_k\beta_k\mu_k$; for a composite
rheology, attribute-dispatched piecewise coefficients built once over the
regions' own splits, so that purely elastic regions (a lithosphere) sit
beside viscoelastic ones with their fixed moduli. An anisotropic rheology,
or a composite with an anisotropic region, returns false. Because the
pointers track the live values, a change of relaxation weights needs only
a reassembly of the blocks, never a re-wiring; the assembly re-reads the
split each time regardless.

\paragraph{Caches.} The primal path keeps the unrelaxed and effective
assemblies in two flavour slots so that stepping schemes that alternate
between them avoid reassembly. Those slots hold primal matrices only: while
the route is on, \code{StashFlavor} declines to stash and
\code{TryRestoreFlavor} refuses to restore (either would clear the dirty
flag over stale blocks), and \code{SetPressureSplitSolver()} invalidates
both slots. Every weight switch therefore reassembles the five forms and
builds a fresh AMG hierarchy; the primal path's preconditioner-reuse
policy does not apply. Under a constant step the weights, and hence the
blocks, are assembled once.

\paragraph{Scope and refusals.} \code{SetPressureSplitSolver()} refuses
(with an \code{MFEM\_VERIFY} naming the reason) unless:
the class supports it (\code{SupportsPressureSplit()}: true for
\code{LinearQuasiStaticClampedProblem} only --- the traction class needs
its rigid-mode projection, the gravitating classes carry their own block
solves); no fluid treatment is set (and \code{SetFluid} refuses while the
route is on); the stiffness has an isotropic split; the rheology carries
no equilibrium mapping (non-natural reference states); the displacement
order is at least 2 (Taylor--Hood); and there is at least one essential
dof, counted across ranks (the route has no rigid-mode projection).

%==========================================================================
\section{Verification}
\label{sec:verification}
%==========================================================================

\code{tests/TestPressureSplit.cpp} runs every gate in 2-D and 3-D, on
simplices and tensor-product elements, at order 2.
\begin{itemize}
\item \emph{Constant strain, exact.} With $\vu = G\vx$ prescribed on the
  whole boundary the Taylor--Hood pair represents the state exactly; the
  displacement is reproduced to $10^{-9}$ and the pressure equals
  $\kappa\tr G$ to $10^{-9}$ relative. This gates the block system, the
  boundary lift and the sign of $p$; the pressure here is entirely the
  constant mode of Remark~\ref{rem:clamped}.
\item \emph{Agreement with the primal solve under refinement.} A beam
  clamped at its base under a surface pull: the two discretisations differ
  by less than 5\% (maximum norm, relative) on the coarse mesh, and the
  gap shrinks by at least a factor $0.7$ under one uniform refinement.
\item \emph{Modulus switches and high contrast.} A two-layer Maxwell
  composite: agreement with the primal solve to 5\% unrelaxed and to 10\%
  with weights $(1,0.3)$ (the tolerance is the interface pressure smearing
  of \S\ref{sec:jump} on a deliberately coarse mesh). With weights
  $(1,10^{-6})$ the primal discretisation locks, so the gate is
  convergence only (and the iteration counts are printed). Clearing the
  weights reproduces the never-weighted solve, and re-setting them the
  weighted one, to $10^{-5}$: the gate against stale blocks.
\item \emph{Elastic beside viscoelastic.} A Maxwell region under a purely
  elastic lid through the composite's piecewise split: agreement with the
  primal solve to 10\% unrelaxed and 20\% at weight $0.3$ (a sharper
  interface pressure jump), and clearing returns the unrelaxed solve to
  $10^{-5}$.
\item \emph{Viscoelastic stepping, exact.} A standard linear solid under
  prescribed uniform strain: the displacement is $\vu = G\vx$ at every
  time, so four backward-Euler steps through the split route --- with the
  weight switches of \code{ImplicitSolve} riding the reassembly --- must
  return it to $10^{-9}$.
\item \emph{Refusals.} The traction class, order 1, and an anisotropic
  rheology are refused with the documented messages.
\end{itemize}
The benchmark driver of the box family exposes the route as \code{-up}.

%==========================================================================
\section{Measured performance and outlook}
\label{sec:results}
%==========================================================================

The \code{contrast} and \code{gradient} models of the viscoelastic box
family (\code{doc/benchmarks.tex}, ``Box: viscosity contrasts''), on 4
ranks, order 2, exponential trapezoid at $\Delta t = 0.25$, 40 solves per
run; iterations are per solve, wall times for the whole run.
\begin{center}
\small
\begin{tabular}{@{}lrrrrr@{}}
\toprule
case & $n_z$ & primal its & split its & primal wall & split wall \\
\midrule
$C = 1$, column & 10 & 55 & 129 & 0.57\,s & 1.34\,s \\
$C = 10^4$, column & 10 & 140 & 203 & 1.7\,s & 2.0\,s \\
$C = 10^4$, slab & 10 & 267 & 424 & 5.6\,s & 7.3\,s \\
$C = 10^6$, column & 10 & 210 & 202 & 2.0\,s & 2.4\,s \\
$C = 10^6$, slab & 10 & 287 & 365 & 3.8\,s & 5.3\,s \\
gradient, $C = 10^6$ & 20 & 1102 & 367 & 7.7\,s & 3.1\,s \\
gradient, $C = 10^6$ & 40 & 2337 & 760 & 21.8\,s & 7.8\,s \\
gradient, $C = 10^6$ & 80 & 4115 & 1756 & 41.1\,s & 19.8\,s \\
\bottomrule
\end{tabular}
\end{center}

\paragraph{Reading.}
\begin{itemize}
\item \emph{Parameter-robust.} The split's column count is flat across
  the contrast ladder, $129\to203\to202$ from $C = 1$ to $10^6$, as
  Proposition~\ref{prop:spectrum} predicts, while the primal count
  quadruples.
\item \emph{A saddle overhead of about $2\times$.} Where nothing is
  ill-conditioned ($C = 1$) the split costs $2.3\times$ the primal solve in
  iterations and wall time: the indefinite spectrum and the extra work per
  iteration of \S\ref{sec:prec}.
\item \emph{Layered versus graded.} On the layered cases the weak layer is
  thin and held between stiff ones, the primal AMG copes, and the primal
  route stays ahead. On the smoothly graded column, where the collapsed
  region is the bulk of the layer and grows in dof count under refinement,
  the primal route collapses and the split wins by $2.5$--$3\times$ in
  iterations and wall time; the primal count reaches 4115 per solve at
  $n_z = 80$. The split's own count also grows under refinement of the
  graded column (roughly doubling per refinement): the variable-coefficient
  and element-shape dependences of \S\ref{sec:inexact}, not a return of
  $\kappa/\mu$.
\item \emph{Locking.} In the harsher two-region configuration of the
  tests (weak-branch weight $10^{-6}$, no stiff support) the gap is
  5--7$\times$ (168 against 932 iterations on one rank, and flat in the
  rank count to four), and there the primal route does not merely slow
  down: it locks, and the two discretisations disagree at $O(1)$
  (\S\ref{sec:jump}). At such contrasts the split is the correct
  discretisation, not only the cheaper solve.
\end{itemize}

\paragraph{Outlook.} The layered box cases are not the production shape.
The regimes the route is aimed at are: graded and laterally varying 3-D
mantle viscosity, the regime of the gradient column; the high-viscosity
lithosphere, effectively elastic, whose base is a rapid but smooth
transition of $\tau$ spanning several decades, so that a range of step
sizes places a steep smooth gradient of $\mu$ there; multi-timescale
rheologies, whose effective operator collapses partially and differently
at each step size (\S\ref{sec:regimes}); and above all the Maxwell
fluid-core treatment, which drives the core's effective shear
$\eps\mu_c\to0$ by design and so produces the most extreme $\kappa/\mu$
ratio in the library (\S\ref{sec:smallmu}). The planned extension is
therefore to the displacement block of the gravitating $(\vu,\phi)$
solve of the mixed self-gravitating class, with the fluid region's
effective shear as one more piecewise region of the split: the potential
block and the bordered rows are unchanged, and the relabelling kernel of
\eqref{eq:kernel} remains the continuation's business.

%==========================================================================
\begin{thebibliography}{99}

\bibitem{Babuska1971}
I.~Babuška,
\emph{Error-bounds for finite element method},
Numer. Math. 16 (1971), 322--333.

\bibitem{BabuskaSuri1992}
I.~Babuška, M.~Suri,
\emph{Locking effects in the finite element approximation of elasticity
problems},
Numer. Math. 62 (1992), 439--463.

\bibitem{BBF}
D.~Boffi, F.~Brezzi, M.~Fortin,
\emph{Mixed Finite Element Methods and Applications},
Springer Series in Computational Mathematics 44, Springer, 2013.
(Inf-sup theory, perturbed saddle-point problems, Taylor--Hood stability
in \S8.8.)

\bibitem{Boffi1997}
D.~Boffi,
\emph{Three-dimensional finite element methods for the Stokes problem},
SIAM J. Numer. Anal. 34 (1997), 664--670.

\bibitem{Braess}
D.~Braess,
\emph{Finite Elements: Theory, Fast Solvers, and Applications in Solid
Mechanics}, 3rd ed., Cambridge University Press, 2007.
(Saddle-point problems with penalty term, Ch.~III \S4; Korn's
inequality and locking, Ch.~VI.)

\bibitem{BrennerScott}
S.~C. Brenner, L.~R. Scott,
\emph{The Mathematical Theory of Finite Element Methods}, 3rd ed.,
Texts in Applied Mathematics 15, Springer, 2008.
(Korn's inequality and volumetric locking, Ch.~11.)

\bibitem{Dain2006}
S.~Dain,
\emph{Generalized Korn's inequality and conformal Killing vectors},
Calc. Var. Partial Differential Equations 25 (2006), 535--540.

\bibitem{ESW}
H.~C. Elman, D.~J. Silvester, A.~J. Wathen,
\emph{Finite Elements and Fast Iterative Solvers}, 2nd ed.,
Oxford University Press, 2014.
(MINRES convergence bounds and block preconditioning of Stokes systems.)

\bibitem{GiraultRaviart}
V.~Girault, P.-A. Raviart,
\emph{Finite Element Methods for Navier--Stokes Equations},
Springer, 1986.

\bibitem{Greenbaum1997}
A.~Greenbaum,
\emph{Iterative Methods for Solving Linear Systems}, SIAM, 1997.

\bibitem{GrinevichOlshanskii2009}
P.~P. Grinevich, M.~A. Olshanskii,
\emph{An iterative method for the Stokes-type problem with variable
viscosity},
SIAM J. Sci. Comput. 31 (2009), 3959--3978.

\bibitem{HensonYang2002}
V.~E. Henson, U.~M. Yang,
\emph{BoomerAMG: a parallel algebraic multigrid solver and
preconditioner},
Appl. Numer. Math. 41 (2002), 155--177.

\bibitem{KarakashianPascal2003}
O.~A. Karakashian, F.~Pascal,
\emph{A posteriori error estimates for a discontinuous Galerkin
approximation of second-order elliptic problems},
SIAM J. Numer. Anal. 41 (2003), 2374--2399.

\bibitem{Klawonn1998}
A.~Klawonn,
\emph{An optimal preconditioner for a class of saddle point problems
with a penalty term},
SIAM J. Sci. Comput. 19 (1998), 540--552.

\bibitem{LWXZ2007}
Y.-J. Lee, J.~Wu, J.~Xu, L.~Zikatanov,
\emph{Robust subspace correction methods for nearly singular systems},
Math. Models Methods Appl. Sci. 17 (2007), 1937--1963.

\bibitem{MalkusHughes1978}
D.~S. Malkus, T.~J.~R. Hughes,
\emph{Mixed finite element methods --- reduced and selective integration
techniques: a unification of concepts},
Comput. Methods Appl. Mech. Engrg. 15 (1978), 63--81.

\bibitem{MardalWinther2011}
K.-A. Mardal, R.~Winther,
\emph{Preconditioning discretizations of systems of partial differential
equations},
Numer. Linear Algebra Appl. 18 (2011), 1--40.

\bibitem{MurphyGolubWathen2000}
M.~F. Murphy, G.~H. Golub, A.~J. Wathen,
\emph{A note on preconditioning for indefinite linear systems},
SIAM J. Sci. Comput. 21 (2000), 1969--1972.

\bibitem{OlshanskiiReusken2006}
M.~A. Olshanskii, A.~Reusken,
\emph{Analysis of a Stokes interface problem},
Numer. Math. 103 (2006), 129--149.

\bibitem{PaigeSaunders1975}
C.~C. Paige, M.~A. Saunders,
\emph{Solution of sparse indefinite systems of linear equations},
SIAM J. Numer. Anal. 12 (1975), 617--629.

\bibitem{RustenWinther1992}
T.~Rusten, R.~Winther,
\emph{A preconditioned iterative method for saddlepoint problems},
SIAM J. Matrix Anal. Appl. 13 (1992), 887--904.

\bibitem{Schoberl1999}
J.~Schöberl,
\emph{Multigrid methods for a parameter dependent problem in primal
variables},
Numer. Math. 84 (1999), 97--119.

\bibitem{ScottVogelius1985}
L.~R. Scott, M.~Vogelius,
\emph{Norm estimates for a maximal right inverse of the divergence
operator in spaces of piecewise polynomials},
RAIRO Modél. Math. Anal. Numér. 19 (1985), 111--143.

\bibitem{SilvesterWathen1994}
D.~Silvester, A.~Wathen,
\emph{Fast iterative solution of stabilised Stokes systems.\ Part II:
using general block preconditioners},
SIAM J. Numer. Anal. 31 (1994), 1352--1367.

\bibitem{Stenberg1984}
R.~Stenberg,
\emph{Analysis of mixed finite element methods for the Stokes problem: a
unified approach},
Math. Comp. 42 (1984), 9--23.

\bibitem{TaylorHood}
C.~Taylor, P.~Hood,
\emph{A numerical solution of the Navier--Stokes equations using the
finite element technique},
Comput. Fluids 1 (1973), 73--100.

\bibitem{Wathen1987}
A.~J. Wathen,
\emph{Realistic eigenvalue bounds for the Galerkin mass matrix},
IMA J. Numer. Anal. 7 (1987), 449--457.

\end{thebibliography}

\end{document}
